\section{Conditional bridges for the unchanged exact return event}
\label{sec:v41-bridges}
The scalar arithmetic coefficient can change with the radius. A path
which is tied to its own endpoint has a more stable conditional limit:
the drift of each long-lived spectral branch cancels before the Fourier
inverse is taken. We prove this cancellation, a conditional fourth-moment
bound, and the passage from the collision clock to the actual return
clock. Conditioning always uses the original event.

For a positive definite matrix $C$, let $\mathbb B_C$ denote the
centered continuous Gaussian bridge with covariance
\[
 \mathbb E\mathbb B_C(s)\mathbb B_C(t)^{\mathsf T}
                         =(s\wedge t-st)C.
\]
Put $\mathsf S_{j,R}=\sum_{a<j}\mathsf f_R\circ T_R^a$, using the
coordinate order in Section~\ref{sec:v41-transition}. Define polygonal
processes by their values at grid points:
\begin{align}
 \mathcal B_{m,R}(j/m)
    &=m^{-1/2}\left(\mathsf S_{j,R}-\frac jm\mathsf S_{m,R}\right),
                                      \label{eq:v41-collision-bridge}\\
 \mathcal Y_{n,R}(l/n)
    &=n^{-1/2}\left(J_{l,R}-\frac ln J_{n,R}\right).
                                      \label{eq:v41-return-bridge}
\end{align}
Both processes are tied to their actual, not nominal, terminal record.
Fix a bounded interval $J$ with $|J|>0$ and write
\begin{equation}\label{eq:v41-exact-conditioning}
 E=E_{n,m,R}(k,t;J)
   =\{K_{n,R}=k,\ N_{n,R}=m,\ T_{n,R}-t\in J\},
 \qquad p_E=\nu_R^*(E).
\end{equation}

\subsection{Pinned characteristic functions before division by the mass}
Let $0=s_0<s_1<\cdots<s_h=1$ be fixed. Write
$m_i=\lfloor ms_i\rfloor-\lfloor ms_{i-1}\rfloor$ and
$t_i^{(m)}=m_i/m$. Given fixed vectors $w_i$, put
$v_i^{(m)}=w_i-\sum_jt_j^{(m)}w_j$. Then
$\sum_i m_i v_i^{(m)}=0$. The exponential of
\[
 m^{-1/2}\sum_i v_i^{(m)}\cdot
    (\mathsf S_{\lfloor ms_i\rfloor,R}
                  -\mathsf S_{\lfloor ms_{i-1}\rfloor,R})
\]
is exactly the characteristic test of the corresponding increments
of \eqref{eq:v41-collision-bridge}. These tests determine its
finite-dimensional laws. Let $N^q_{n,m,R}$ be the expectation of this
exponential times the exact displacement, occupation and two-section
indicators, with a roof test $q(S_{m,R}-t)$, under $c^{-1}\nu$.

\begin{lemma}[Uniform pinned factorization on a fixed roof interval]
\label{lem:v41-pinned-factorization}
For the preceding fixed observation times and vectors, uniformly in
$R,n,k,t$,
\begin{align}
 m^2\left[N^J_{n,m,R}(k,t)
     -\Phi_{m,R}\,P_{n,m,R}(k,t;J)\right]&\longrightarrow0,
                                      \label{eq:v41-pinned-factorization}\\
 \Phi_{m,R}
   &=\exp\left(-\frac12\sum_i t_i^{(m)}
           (v_i^{(m)})^{\mathsf T}\Omega_Rv_i^{(m)}\right).
                                      \label{eq:v41-bridge-characteristic}
\end{align}
Here $N^J$ denotes the test with $\one_J$. The factor in
\eqref{eq:v41-bridge-characteristic} tends to the characteristic
function of the corresponding increments of $\mathbb B_{\Omega_R}$.
No lower bound for the conditioning probability is used in this lemma.
\end{lemma}
\begin{proof}
First take an integrable band-limited roof test. Fourier inversion gives
the chronological product
\begin{equation}\label{eq:v41-pinned-product}
 \ell\left(M_{\eta_R}
  \mathscr T_{R,\theta+v_h^{(m)}/\sqrt m}^{m_h}\cdots
  \mathscr T_{R,\theta+v_1^{(m)}/\sqrt m}^{m_1}(\eta_R\nu)\right)
\end{equation}
inside the same inverse as in \eqref{eq:v40-fixed-spectral-inverse}.
All block lengths are positive multiples of $m$ for large $m$.
On the compact nonresonant complement, the small frequency shifts
remain in common decaying neighborhoods, so this part is exponentially
small. Near a peak $a_j$, set $\theta=a_j+y/\sqrt m$ and use the local
splitting on every block. A term with a complementary power decays
exponentially. The product of the principal projections has endpoint
amplitude $A_j(R,a_j)+O((1+|y|)/\sqrt m)$: at the common center the
projection is idempotent and of rank one.

The logarithm of the principal eigenvalue product is
\begin{align*}
 m\log z_j+i\sqrt m\,\mu_j\cdot y
   -\frac12\sum_i t_i^{(m)}(y+v_i^{(m)})^{\mathsf T}
                       H_j(y+v_i^{(m)})
                 +O(m^{-1/2}(1+|y|)^3).
\end{align*}
The drift contribution of the $v_i^{(m)}$ is zero. Their cross term
with $y$ is zero for the same exact reason. The Gaussian integral on
this chart is consequently the unmarked term from
\eqref{eq:v41-test-peak-sum}, multiplied by
\[
 \exp\left(-\tfrac12\sum_i t_i^{(m)}
                      (v_i^{(m)})^{\mathsf T}H_jv_i^{(m)}\right).
\]
The real quadratic bounds dominate the product, since
$\sum_i t_i^{(m)}|y+v_i^{(m)}|^2
=|y|^2+\sum_i t_i^{(m)}|v_i^{(m)}|^2$.
The rescaling and tail argument of
Theorem~\ref{thm:v41-uniform-transition} therefore applies uniformly.

On a compact parameter support, $\kappa_j=0$ implies $H_j=\Omega_R$.
It follows that
\begin{equation}\label{eq:v41-damped-curvature-continuity}
 \sup_R e^{-m\kappa_j(R)}\|H_j(R)-\Omega_R\|\longrightarrow0.
\end{equation}
To see this, use continuity near the zero set of $\kappa_j$, and a
positive minimum of $\kappa_j$ on the compact complement of any such
neighborhood. All Gaussian integrals and other scalar factors are
bounded. Thus replacing every $H_j$ in the extra characteristic factor
by $\Omega_R$ changes the total by $o(1)$. This proves the band-limited
version of \eqref{eq:v41-pinned-factorization}, including parameter
sequences which approach a change of arithmetic type.

For the sharp interval take its upper envelope $q_B^+$ and lower
envelope $q_B^-$. Since the inserted path exponential has modulus one,
\[
 |N^J-N^{q_B^+}|\le P^{q_B^+}-P^J
                         \le P^{q_B^+}-P^{q_B^-}.
\]
The uniform scalar formula bounds the limsup of $m^2$ times the last
difference by $C/B$. The same estimate pays the change in the unmarked
probability. Fix $B$, let $m\to\infty$, and then let $B\to\infty$.
This proves \eqref{eq:v41-pinned-factorization}. Finally the displayed
quadratic form is the elementary covariance formula for independent
Gaussian increments tied to their sum.
\end{proof}

\subsection{Fourth moments under a microscopic constraint}
\begin{lemma}[A derivative estimate centered at a moving spectral peak]
\label{lem:v41-peak-derivatives}
On each peak chart, for a real unit vector $e$, $0\le p\le4$ and
$L\ge1$, there are uniform $C,a>0$ and $\rho<1$ such that
\begin{align}
 \left\|\left.\frac{d^p}{dz^p}
    \left(e^{-iz e\cdot\mu_j(R)}
                  \mathscr T_{R,\theta+ze}\right)^L
                         \right|_{z=0}\right\|
 \le C L^{p/2}\left(e^{-aL|\theta-a_j(R)|^2}+\rho^L\right).
                                      \label{eq:v41-peak-derivatives}
\end{align}
For $L=0$ the derivative is zero when $p>0$, and is the identity when
$p=0$. On a compact nonresonant part of a fixed band, the corresponding
uncentered derivatives have a bound $C_p\rho^L$.
\end{lemma}
\begin{proof}
Put $x=\theta-a_j$. The local analytic splitting is uniform in a
complex neighborhood. On a circle $|z|=d/\sqrt L$, with $d$ fixed
small, its principal eigenvalue after centering satisfies
\[
 \left|e^{-iz e\cdot\mu_j}\lambda_j(R,\theta+ze)\right|^L
  \le e^{-L\kappa_j}
       \exp\{-a_0L|x|^2+C L|x||z|+C L|z|^2\}
  \le C e^{-a_1L|x|^2}.
\]
The first inequality follows from the zero real gradient at the peak,
the positive real Hessian there and uniform third derivatives, reducing
the chart once. Completing the square gives the second. The projection
has uniformly bounded norm on this circle. Its complementary power is
bounded by $C\rho_0^L e^{C\sqrt L}$, which is at most $C\rho^L$ for
some $\rho_0<\rho<1$. Cauchy's derivative estimate supplies the factor
$L^{p/2}$. Increasing $C$ treats the finite initial range of $L$.
On the nonresonant complement use a fixed complex circle small enough
that the spectral contour still has radius less than one. Cauchy's
estimate on that circle proves the last assertion.
\end{proof}

\begin{proposition}[An unnormalized pinned fourth moment]
\label{prop:v41-pinned-fourth}
For $0\le r<r+l\le m$, $l\ge1$, uniformly in all radii and targets,
\begin{equation}\label{eq:v41-pinned-fourth}
 \int_E\left|\mathsf S_{r+l,R}-\mathsf S_{r,R}
                         -\frac lm\mathsf S_{m,R}\right|^4d\nu_R^*
                              \le C_J l^2m^{-2}.
\end{equation}
The numerator uses the original exact return event
\eqref{eq:v41-exact-conditioning}. The estimate has not yet been
divided by its probability.
\end{proposition}
\begin{proof}
It suffices to prove the scalar estimate in each coordinate. Majorize
$\one_J$ by one fixed nonnegative band-limited upper envelope. The
other source factors and the fourth power are nonnegative, so this
increases the integral. Put $\alpha=l/m$ and $s=m-r-l$. The scalar
fourth moment inside the Fourier inverse is the fourth derivative at
$z=0$ of
\[
 \ell\left(M_{\eta_R}
  \mathscr T_{R,\theta-\alpha ze}^{s}
  \mathscr T_{R,\theta+(1-\alpha)ze}^{l}
  \mathscr T_{R,\theta-\alpha ze}^{r}(\eta_R\nu)\right).
\]
The identity follows by differentiating the exact physical pairing.
Every one-collision observable is bounded, so differentiation under
the finite-time integral and the fixed-band inverse is legitimate.

On peak chart $j$, center each of the three operators at $\mu_j(R)$.
The scalar factors introduced by this centering multiply to one,
since $-\alpha(s+r)+(1-\alpha)l=0$. Apply
Lemma~\ref{lem:v41-peak-derivatives} to each differentiated factor.
The sum of the size factors in Leibniz's formula is bounded by
\[
 C\{\alpha(\sqrt r+\sqrt s)+(1-\alpha)\sqrt l\}^4
       \le C l^2.
\]
Indeed $\alpha(\sqrt r+\sqrt s)\le\sqrt2\,l/\sqrt m
\le\sqrt{2l}$. At least one block has length at least $m/3$.
Its Gaussian factor has integral $O(m^{-2})$ in four frequencies;
its complementary factor is exponentially small. The other factors
are bounded by the preceding size estimate. Thus the integral over
each chart is at most $C_Jl^2m^{-2}$. A zero-length factor is interpreted
as in the lemma and causes no exception.

On the compact nonresonant part, derivatives of the longest block
decay exponentially. The remaining fixed-order derivatives can be
bounded polynomially in $m$, either by the lemma or by differentiating
a product of uniformly bounded real-frequency powers. The resulting
$C_J(1+m)^4\rho^{m/3}$ is at most $C_Jl^2m^{-2}$ after increasing
the constant. Multiplication by the integrable Fourier transform of
the fixed envelope and summation over the finite cover prove the
scalar estimate. Summing coordinates proves the vector assertion.
The fixed factor $c^{-1}$ is absorbed in $C_J$.
\end{proof}

\begin{theorem}[A uniform conditional collision bridge]
\label{thm:v41-collision-conditional-bridge}
Fix $M,\epsilon>0$ and $J$ as above. Uniformly over the original events
with
\begin{equation}\label{eq:v41-denominator-class}
 |Z|\le M,\qquad m^2p_E\ge\epsilon,
\end{equation}
one has, on $C([0,1],\R^4)$ with its supremum norm,
\begin{equation}\label{eq:v41-collision-conditional-bridge}
 d_{\rm BL}\left(\operatorname{Law}_{\nu_R^*(\cdot\mid E)}
                        (\mathcal B_{m,R}),
                         \operatorname{Law}(\mathbb B_{\Omega_R})\right)
                        \longrightarrow0.
\end{equation}
\end{theorem}
\begin{proof}
Divide \eqref{eq:v41-pinned-factorization} by $p_E$, using
\eqref{eq:v41-denominator-class}. This proves uniform convergence of
finite-dimensional characteristic functions to those of the Gaussian
bridge. Polygonal interpolation has no effect on these limits, since
the collision record is bounded and the total centered endpoint is
bounded on the stated event class.

By Proposition~\ref{prop:v41-pinned-fourth}, a grid increment of the
normalized process has conditional fourth moment at most
$C_{J,\epsilon}(l/m)^2$. Linear interpolation gives the same bound
$C_{J,\epsilon}|t-s|^2$ at arbitrary times: within one grid interval
its fourth moment is multiplied by the fourth power of the fractional
length; an arbitrary increment is the sum of two such end pieces
and a complete grid interval. The fourth-moment continuity criterion
therefore gives uniform tightness in continuous path space.

For the bounded-Lipschitz formulation, take any hypothetical violating
sequence. Compactness gives a convergent subsequence of its radii.
Uniform tightness and the finite-dimensional limits identify every
subsequential law with the Gaussian bridge at that limiting radius.
Continuity of $\Omega_R$ gives continuity of these bridge laws, for
example by coupling with $\Omega_R^{1/2}$ times one standard bridge.
This contradicts the violating sequence. The proof uses no mixing
assumption for the induced return map.
\end{proof}

\subsection{Returning to the genuine return clock}
\begin{theorem}[The microscopic bridge of the actual return record]
\label{thm:v41-actual-return-bridge}
Under the same event class \eqref{eq:v41-denominator-class}, uniformly
in the radius and central targets,
\begin{equation}\label{eq:v41-actual-return-bridge}
 d_{\rm BL}\left(\operatorname{Law}_{\nu_R^*(\cdot\mid E)}
                         (\mathcal Y_{n,R}),
                         \operatorname{Law}(\mathbb B_{D_R})\right)
                         \longrightarrow0\qquad(m\to\infty).
\end{equation}
At every fixed radius the conclusion holds uniformly on central target
compact sets over all positive arithmetic residue classes in
Theorem~\ref{thm:v40-arithmetic-return-LLT}; no additional denominator
hypothesis is then needed. The conditioning fixes both lattice
displacement coordinates and the collision count exactly, and keeps
the roof in the specified fixed interval.
\end{theorem}
\begin{proof}
Let $\mathcal X_{m,R}$ be the polygonal path with values
$m^{-1/2}\sum_{a<j}X_R(T_R^ax)$ at $j/m$. On $E$ its endpoint equals
$Z+O_J(m^{-1/2})$. It is consequently tight under the conditional laws,
because $\mathcal X_{m,R}(s)=\mathcal B_{m,R}(s)
+s\mathcal X_{m,R}(1)$ and Theorem~\ref{thm:v41-collision-conditional-bridge}
applies. In particular
\[
 \sup_{j\le m}|A_{j,R}-cj|/m\longrightarrow0
\]
in conditional probability, uniformly in the event class. At a genuine
return, the half-open count identity is $A_{N_{l,R},R}=l$.
On $E$ every $N_{l,R}$ with $0\le l\le n$ lies in $[0,m]$ and
$N_{n,R}=m$. Since $n/m=c+O(m^{-1/2})$, the preceding bound implies
\begin{equation}\label{eq:v41-conditional-return-clock}
 \max_{0\le l\le n}\left|\frac{N_{l,R}}m-\frac ln\right|
                              \longrightarrow0
\end{equation}
in those same conditional probabilities. No comparison with a
substitute conditioning event has been made.

Let $L_R$ be the invertible matrix in
\eqref{eq:v37-covariance-change}. The exact compensation and
$X_R=L_Rh_R$ give
\[
 J_{l,R}-l\bar G_R
     =L_R^{-1}\sum_{a<N_{l,R}}X_R(T_R^ax).
\]
Writing $\theta_l=N_{l,R}/m$, subtraction of the terminal record gives
at every return grid point the exact identity
\begin{align}\label{eq:v41-pinned-time-change}
 \mathcal Y_{n,R}(l/n)
   =\sqrt{m/n}\,L_R^{-1}
       \left\{\mathcal B_{m,R}(\theta_l)
              +(\theta_l-l/n)\mathcal X_{m,R}(1)\right\}.
\end{align}
The second term tends to zero in supremum norm by
\eqref{eq:v41-conditional-return-clock} and the bounded endpoint.
Uniform tightness with continuous limits permits replacing
$\mathcal B_{m,R}(\theta_l)$ by $\mathcal B_{m,R}(l/n)$.
Linear interpolation between the return grid points preserves this
conclusion; it follows from the same modulus of continuity and does
not require a separate tail estimate divided by a rare-event mass.
Finally $m/n\to c^{-1}$ and
\[
 c^{-1}L_R^{-1}\Omega_R(L_R^{-1})^{\mathsf T}=D_R.
\]
This proves \eqref{eq:v41-actual-return-bridge} and its uniformity.
For fixed $R$, the finitely many positive arithmetic coefficients
have a positive minimum. The sharp-interval local theorem supplies
\eqref{eq:v41-denominator-class} with a fixed $\epsilon_{R,M,J}>0$,
which proves the last assertion.
\end{proof}

\paragraph{Topology and the remaining raw problem.}
These are continuous-path limits under the unchanged exact discrete
return constraints and a fixed positive roof interval. The lower mass
condition is explicit in the radius-uniform assertion, and is proved
on every positive fixed-radius arithmetic class. A pointwise roof
conditioning law, arbitrary selected-path Gaussian amplitudes, and the
common raw density correction do not follow by division by the interval
length. The original pointwise theorem retains those requirements.
